%
% 23-anwendungen.tex -- Anwendungen
%
% (c) 2026 Prof Dr Andreas Müller
%
\subsection{Anwendungen
\label{buch:holomorph:geometrie:subsection:anwendungen}}
Konforme Abbildungen ermöglichen, komplizierte Gebiete auf 
komplex differenzierbare Art auf sehr einfache Gebiete abzubilden.
Zum Beispiel kann man nach einer harmonischen Funktion mit 
gegebenen Randwerten auf dem Gebiet fragen.
Die konforme Abbildung macht daraus das Problem, eine harmonische
Funktion auf einem Kreisgebiet mit gegebenen Randwerten zu finden.
Für Lösung dieses Problem gibt es aber sogar eine Formel von Poisson.
Durch Zusammensetzen der Lösung von Poisson mit der konformen
Abbildung entsteht eine Lösung des schwierigen Problems auf dem
ursprünglichen Gebiet.

Diese Idee, durch Abbildung des Gebietes eine Aufgabenstellung
auf eine bereits gelöste Aufgabe zu transformieren, wird in
zwei weiteren Kapiteln dieses Seminars verwendet.
In Kapitel~\ref{chapter:elektro} wird die Elektrostatik in der
Ebene untersucht.
Da in der Elektrostatik alles durch das elektrostatische Potential
gegeben ist, genügt es Potentialfunktionen zu finden.
Für komplexe Geometrie der berandenden Leiter kann dies eine
sehr schwierige Aufgabe sein.
Das Potential eines unregelmässigen Körpers kann aber durch
eine Abbildung der Ebene ausserhalb des Körpers auf das
äussere der Einheitskreisscheibe gefunden werden,
da das Potential eines Kreises $\sim\frac1r$ ist.

Zweidimensionale laminare Strömungen lassen sich durch die sogenannte
Stromfunktion beschreiben, die eine harmonische Funktion ist.
Es ist jedoch im allgemeinen schwierig, diese Funktion für eine
komplexe Geometrie des umströmten Körpers, zum Beispiel eines Flügelprofils,
zu finden.
Für die Umströmung eines Kreises kann man jedoch eine Formel
angeben.
Es muss jetzt nur noch eine Abbildung des Aussenraumes des Flügelprofils
auf den Aussenraum des Kreises gefunden werden.
Diese Idee ermöglicht sogar eine Formel für den Auftrieb eines Flügels
herzuleiten, sie wird in Kapitel~\ref{chapter:joukowski} genauer
ausgeführt.


